% =============================================================================
\section{Discussion and Limitations}\label{sec:discussion}
% Target length: ~0.75 page (~350-400 words)
% =============================================================================

\paragraph{Decision-aware loss carries the forecasting story.}
Ablating the multi-quantile pinball loss back to MSE
(\Cref{tab:forecast_accuracy}, \emph{w/o decision-aware loss})
inflates $h = 14$ RMSE by $29\%$ ($14.46 \to 18.70$) and the
under\-estimation rate at $h = 28$ rises from $10.3\%$ to $44\%+$
across the classical baselines. The level-and-trend anchor and the
PINN parameter features are individually necessary too --- removing
the parameter features inflates $h = 14$ RMSE almost two-fold
($14.46 \to 27.25$). Together these results support the central
methodological claim of the paper: under operational asymmetry,
choosing the loss is as load-bearing as choosing the architecture.

\paragraph{Robust MILP wins the routing game, not the coverage game.}
At the operating budget ($B = 213$~beds, $\tau = 240$~min) every
budget-exhausting policy in \Cref{tab:allocation} reaches the same
expected and worst-case unmet floor; what separates them is transfer
burden. The robust MILP achieves $672$~bed-km, $4.6\times$ below
population-proportional ($3120$) and $1.7\times$ below greedy
shortage-first ($1140$), at no cost in coverage. The deterministic
MILP --- optimising only the median scenario --- spends only $51$ of
the $213$ available surge beds and pays for this with $\sim 3\times$
the expected and worst-case unmet of the budget-exhausting policies,
illustrating the operational value of optimising the
$\pi_s$-weighted expectation rather than the median alone.

\paragraph{The CVaR machinery is informative even when non-binding.}
The $\lambda_3$ sweep (\Cref{tab:e6}) shows that the robust MILP's
allocation coincides with the expectation-minimising one for every
$\lambda_3 \in [0, 4]$ at the operating budget. This is not a defect
but a measurement: at $20\%$ surge budget the high-scenario tail
($q^{(0.9)}$) is fully cleared by the budget-exhausting expectation
solution, so the CVaR constraint is slack. A tighter budget ($\leq
15\%$, where the worst-case scenario admits residual unmet) would
re-engage $\lambda_3$ as a free control; this is the regime in which
the robust formulation would diverge from the deterministic median
optimum even after the LP-slave re-evaluation.

\paragraph{Forecast quality matters at the q$^{0.9}$ margin.}
The five-forecaster E5 sweep (\Cref{tab:e5}) shows that all
forecasters but seasonal-naive($7$) over-state the realised peak at
the chosen origin; PinnGRU is the most conservative ($q^{(0.9)}$ sum
of $1421$~beds vs.\ a realised peak total of $797$). Because the
Delta-peak baseline is generous relative to the realised Omicron MV
peak, realised unmet is zero across the board including the oracle.
The operational reading is that the value of forecast quality in this
regime lies in \emph{calibration of the $q^{(0.9)}$ margin}: a tighter
$q^{(0.9)}$ buys efficiency (fewer pre-positioned beds), but the same
$q^{(0.9)}$ becomes the safety margin in a future wave whose realised
peak exceeds the Delta-peak ceiling. Seasonal-naive's
$q^{(0.9)}$-equivalent of $745$~beds falls $52$~beds below the
realised peak, the only baseline that would have under-provisioned.

\paragraph{Limitations.}
Inter-regional travel times are approximated by detour-scaled
great-circle distances between regional centroids; an OSRM-derived
matrix would refine particular pair-wise routings without changing
qualitative conclusions. The compartmental model assumes well-mixed
regional populations and a fixed clinical-pathway structure; coupling
to age-stratified pathways~\cite{shamseddin2025optimization} is a
natural extension. Costs in the objective are reported as bed and
bed-km quantities; converting to monetary units would require
commissioning-context calibration. Finally, the framework is
configured around NHS England's seven-region structure; transfer to
other health systems would require re-specifying demand units,
capacity categories, and travel infrastructure.
